% ============================================================
%  ANÁLISIS BIBLIOMÉTRICO: TEA & ALTAS CAPACIDADES EN ACADEMIA
%  Para compilar en Overleaf (pdfLaTeX + biber)
% ============================================================
\documentclass[12pt, a4paper]{article}

% ── Codificación y lengua ──────────────────────────────────
\usepackage[utf8]{inputenc}
\usepackage[T1]{fontenc}
\usepackage[spanish, es-tabla]{babel}

% ── Tipografía ─────────────────────────────────────────────
\usepackage{lmodern}
\usepackage{microtype}

% ── Geometría y espaciado ──────────────────────────────────
\usepackage[top=2.5cm, bottom=2.5cm, left=3cm, right=2.5cm]{geometry}
\usepackage{setspace}
\onehalfspacing

% ── Matemáticas ────────────────────────────────────────────
\usepackage{amsmath, amssymb}

% ── Gráficas y tablas ──────────────────────────────────────
\usepackage{graphicx}
\usepackage{booktabs}
\usepackage{longtable}
\usepackage{multirow}
\usepackage{array}
\usepackage{tabularx}
\usepackage{caption}
\usepackage{subcaption}
\usepackage{float}
\usepackage[table, xcdraw]{xcolor}

% ── Hipervínculos ──────────────────────────────────────────
\usepackage[colorlinks=true,
            linkcolor=blue!60!black,
            citecolor=blue!60!black,
            urlcolor=blue!60!black]{hyperref}

% ── Bibliografía (biblatex + biber) ───────────────────────
\usepackage[backend=biber,
            style=apa,
            sortcites=true,
            sorting=nyt,
            maxcitenames=3,
            maxbibnames=10]{biblatex}
\addbibresource{referencias.bib}   % archivo .bib separado

% ── Utilidades ─────────────────────────────────────────────
\usepackage{enumitem}
\usepackage{csquotes}
\usepackage{pdflscape}
\usepackage{rotating}
\usepackage{lipsum}

% ── Colores personalizados ─────────────────────────────────
\definecolor{azuloscuro}{RGB}{46, 74, 122}
\definecolor{rojooscuro}{RGB}{192, 57, 43}
\definecolor{verdoscuro}{RGB}{92, 122, 46}
\definecolor{grisFondo}{RGB}{245, 242, 235}

% ── Cajas de advertencia / nota ───────────────────────────
\usepackage{mdframed}
\newmdenv[backgroundcolor=grisFondo,
          linecolor=azuloscuro,
          linewidth=1pt,
          leftline=true, rightline=false,
          topline=false, bottomline=false,
          innerleftmargin=8pt, innerrightmargin=8pt,
          innertopmargin=6pt, innerbottommargin=6pt]{notabox}

% ── Encabezados ────────────────────────────────────────────
\usepackage{fancyhdr}
\pagestyle{fancy}
\fancyhf{}
\rhead{\small TEA, Altas Capacidades e Investigación}
\lhead{\small Análisis Bibliométrico · Lens.org}
\cfoot{\thepage}
\renewcommand{\headrulewidth}{0.4pt}

% ── Metadatos PDF ──────────────────────────────────────────
\hypersetup{
  pdftitle={Análisis bibliométrico sobre la intersección del Trastorno del Espectro Autista y las Altas Capacidades en investigadores y académicos},
  pdfauthor={},
  pdfkeywords={TEA, altas capacidades, doble excepcionalidad, bibliometría, academia, Colombia, Latinoamérica}
}

% ═══════════════════════════════════════════════════════════
\begin{document}
% ═══════════════════════════════════════════════════════════

% ── PORTADA ───────────────────────────────────────────────
\begin{titlepage}
\centering
\vspace*{1cm}

{\color{azuloscuro}\rule{\linewidth}{2pt}}
\vspace{0.5cm}

{\Large\bfseries Análisis Bibliométrico sobre la Intersección del\\[6pt]
Trastorno del Espectro Autista y las Altas Capacidades\\[6pt]
en Investigadores y Académicos Profesionales}

\vspace{0.4cm}
{\color{rojooscuro}\rule{\linewidth}{0.8pt}}

\vspace{0.8cm}

{\large\itshape Una revisión del estado del arte con perspectiva\\
global, latinoamericana y colombiana}

\vspace{1.2cm}

% Autores — completar
{\normalsize
\textbf{[Autor/es]}\\[4pt]
\textit{[Afiliación institucional]}\\[4pt]
\textit{[Correo electrónico]}
}

\vspace{1.2cm}

{\normalsize
\textbf{Tipo de manuscrito:} Revisión bibliométrica\\[4pt]
\textbf{Fuente de datos:} Lens.org\\[4pt]
\textbf{Fecha de la búsqueda:} Junio de 2026\\[4pt]
\textbf{Período analizado:} 1998--2026
}

\vspace{1.5cm}

\begin{notabox}
\small
\textbf{Palabras clave:} trastorno del espectro autista; altas capacidades; doble
excepcionalidad; \textit{twice-exceptional}; bibliometría; análisis cienciométrico;
academia; investigación; diagnóstico tardío; enmascaramiento; Colombia;
Latinoamérica; neurodiversidad.
\end{notabox}

\vfill
{\color{azuloscuro}\rule{\linewidth}{2pt}}
\end{titlepage}

% ── RESUMEN ───────────────────────────────────────────────
\newpage
\begin{abstract}
\noindent
\textbf{Antecedentes y objetivo.} La co-ocurrencia del Trastorno del Espectro Autista
(TEA) y las Altas Capacidades Intelectuales (AACC) —conocida como doble
excepcionalidad o \textit{twice-exceptionality} (2e)— ha ganado atención investigadora
en la última década, aunque la producción científica sigue siendo escasa y concentrada
casi exclusivamente en poblaciones infantiles y escolares. Los individuos adultos con
perfil 2e que se desempeñan en entornos académicos o científicos representan un
subgrupo particularmente invisible en la literatura. Este trabajo tiene como objetivo
mapear y caracterizar la producción científica global sobre la intersección TEA--AACC
con especial énfasis en investigadores y académicos profesionales, y situar el estado
del campo en los contextos latinoamericano y colombiano.

\noindent
\textbf{Métodos.} Se realizó un análisis bibliométrico a partir de cinco estrategias
de búsqueda booleana ejecutadas en Lens.org en junio de 2026. El corpus final, tras
deduplicación por DOI e identificador Lens, comprende \textbf{3.966 registros} (1834--2026),
de los cuales 2.838 (71,6\,\%) son artículos de revista, distribuidos en 1.635 fuentes
distintas. Se aplicaron indicadores de productividad (Ley de Bradford, Ley de Lotka),
análisis de coautoría, evolución temporal y análisis de co-ocurrencia de términos en
títulos, con visualizaciones generadas en Python (\texttt{bibtexparser}, \texttt{pandas},
\texttt{networkx}, \texttt{wordcloud}).

\noindent
\textbf{Resultados.} La producción científica creció a una tasa anual compuesta del
16,4\,\% entre 2010 y 2020, con un pico de 425--424 registros en 2021--2022. El 78,3\,\%
del corpus total se publicó después de 2015. El núcleo Bradford comprende 113 revistas
que concentran el primer tercio de la producción, lideradas por \textit{Frontiers in
Psychology} ($n=65$), \textit{Journal of Autism and Developmental Disorders} ($n=41$) y
\textit{Focus on Exceptional Children} ($n=33$). La distribución de autoría sigue el
patrón de Lotka (90,9\,\% de los 20.986 autores únicos tienen una sola publicación;
solo 24 autores tienen $\geq10$ publicaciones). Solo el 6,5\,\% del corpus menciona
explícitamente adultos, investigadores o entornos de educación superior en sus
títulos. La literatura en español (343 registros, 8,6\,\% del corpus) presenta un
solapamiento mínimo (5,0\,\%) con el corpus anglosajón, y los registros
colombianos identificados suman únicamente 8 documentos.

\noindent
\textbf{Conclusiones.} El campo presenta un crecimiento exponencial pero con brechas
estructurales críticas: ausencia de investigación sobre adultos 2e en contextos
académicos, subrepresentación de la perspectiva iberoamericana, escasez de estudios
sobre enmascaramiento y diagnóstico tardío en profesionales, e invisibilidad casi
total del contexto colombiano. Se formulan seis líneas prioritarias de investigación
futura.
\end{abstract}

\vspace{0.5cm}
\noindent\textbf{Clasificación JEL / MSC:} I23, I24, J24 \quad
\textbf{Diseño del estudio:} Bibliométrico, transversal retrospectivo

\newpage
\tableofcontents
\newpage

% ═══════════════════════════════════════════════════════════
\section{Introducción}
\label{sec:intro}
% ═══════════════════════════════════════════════════════════

El reconocimiento científico de la doble excepcionalidad como fenómeno distinto y
merecedor de atención específica tiene raíces que se remontan a los trabajos pioneros
de Leta Hollingworth en los años veinte del siglo~XX \parencite{hollingworth1923},
pero su formalización conceptual en la literatura educativa contemporánea no se
consolida hasta la década de 1990. La expresión \textit{twice-exceptional} (2e),
acuñada para describir a individuos que presentan simultáneamente capacidades
intelectuales superiores y alguna condición de neurodesarrollo, neurodiversidad o
discapacidad, fue adoptada con mayor sistematicidad a partir de ese período
\parencite{baum2004}.

Dentro del espectro de condiciones que configuran la doble excepcionalidad, la
co-ocurrencia de Altas Capacidades Intelectuales (AACC) con el Trastorno del Espectro
Autista (TEA) es, paradójicamente, una de las más frecuentes y de las más difíciles
de detectar \parencite{foleynicpon2011}. La razón es múltiple: los perfiles cognitivos
elevados pueden enmascarar las manifestaciones conductuales del TEA \parencite{boschi2016},
los instrumentos de cribado desarrollados para poblaciones neurotípicas infravaloran
las capacidades en personas con TEA \parencite{assouline2012}, y las expectativas
sociales hacia individuos con alto rendimiento académico inhiben la búsqueda de
evaluación diagnóstica \parencite{gelbar2022}.

Esta tensión se agudiza en entornos académicos y científicos. Los profesionales con
TEA y AACC que trabajan en universidades o centros de investigación constituyen un
perfil singularmente invisibilizado: su alta funcionalidad intelectual y la capacidad
de compensar —mediante estrategias de enmascaramiento (\textit{masking} o
\textit{camouflaging})— las dificultades asociadas al TEA permiten que permanezcan
sin diagnóstico durante décadas \parencite{mandy2019}. Las consecuencias de esta
invisibilidad son bien documentadas a nivel clínico —agotamiento neurológico,
burnout, ansiedad crónica— pero permanecen casi sin explorar en su dimensión
específicamente académica y profesional \parencite{price2022}.

Desde la perspectiva latinoamericana, la situación es aún más precaria. La
investigación sobre doble excepcionalidad en la región es embrionaria, la mayoría
de los sistemas educativos nacionales carecen de políticas diferenciadas para esta
población, y los estudios disponibles sobre adultos con TEA son escasos y, en el
caso colombiano, prácticamente inexistentes \parencite{conejeros2018}.

El presente trabajo tiene como objetivo central caracterizar bibliométricamente el
estado actual del conocimiento científico sobre la intersección TEA--AACC, con
énfasis particular en individuos adultos que se desempeñan en entornos de
investigación y academia, y con un análisis específico de la producción
iberoamericana y colombiana. Para ello, se aplican técnicas bibliométricas
estándar sobre un corpus de 3.966 registros recuperados de Lens.org.

\subsection{Justificación del estudio}

Tres razones fundamentan la oportunidad de este análisis. Primera, la \textbf{brecha
empírica}: los estudios bibliométricos previos sobre el campo 2e \parencite{luor2021,
gelbar2022} identificaron corpora de apenas 32--95 artículos, lo que sugiere una
cobertura de bases de datos restringida (principalmente WoS/SSCI) y unos criterios de
inclusión centrados en educación primaria y secundaria. El uso de Lens.org, que
agrega más de 200 millones de registros de PubMed, Crossref y OpenAlex, permite
ampliar y actualizar este panorama. Segunda, la \textbf{brecha temporal}: ningún
estudio bibliométrico publicado cubre el período 2021--2026, que nuestro corpus
revela como el de mayor producción histórica. Tercera, la \textbf{brecha geográfica}:
la perspectiva latinoamericana y colombiana específicamente permanece ausente de los
análisis existentes.

% ═══════════════════════════════════════════════════════════
\section{Marco conceptual}
\label{sec:marco}
% ═══════════════════════════════════════════════════════════

\subsection{Trastorno del Espectro Autista: definición y prevalencia}

El DSM-5 \parencite{apa2013} define el TEA como un trastorno del neurodesarrollo
caracterizado por: (a) déficits persistentes en la comunicación e interacción social
en múltiples contextos, y (b) patrones de comportamiento, intereses o actividades
restrictivos y repetitivos. La adopción del modelo dimensional en el DSM-5,
eliminando las subcategorías previas (síndrome de Asperger, trastorno generalizado
del desarrollo no especificado), implicó un cambio sustantivo que amplió el rango
diagnóstico y reconoció la heterogeneidad del espectro.

Las estimaciones globales de prevalencia han aumentado sostenidamente: de 1 en 166
niños en 2004 \parencite{cdc2004} a 1 en 36 en 2023 \parencite{maenner2023}. Este
incremento refleja tanto una mayor conciencia diagnóstica como el impacto de las
reformas nosológicas. Históricamente, se estimaba que el 70--75\,\% de las personas
con TEA presentaban discapacidad intelectual \parencite{russell2019}; en la actualidad,
esta cifra se ha invertido, situándose en torno al 33\,\% \parencite{zeidan2022},
lo que subraya la prevalencia de perfiles TEA sin discapacidad intelectual.

\subsection{Altas Capacidades: modelos y conceptualización}

El concepto de altas capacidades (AACC) carece de una definición universal, lo que
genera dificultades para la comparación entre estudios. Los modelos más influyentes
incluyen: la definición psicométrica clásica (CI$\geq$130, aproximadamente 2 desviaciones
estándar por encima de la media \parencite{terman1925}), el modelo triádico de
Renzulli \parencite{renzulli1978} (que combina habilidad por encima del promedio,
creatividad y compromiso con la tarea), y el modelo de Munich \parencite{heller2001},
que reconoce dominios múltiples de talento.

En el contexto iberoamericano, el término \textit{altas capacidades intelectuales}
tiende a englobar tanto la superdotación (alto potencial en múltiples dominios) como
el talento específico (excelencia en un dominio determinado) y el prodigio
\parencite{gagné2004}. La diversidad terminológica —\textit{superdotación},
\textit{sobredotación}, \textit{altas habilidades/superdotação} en el contexto
brasileño— dificulta la sistematización bibliométrica en castellano y portugués.

\subsection{Doble excepcionalidad: la intersección TEA--AACC}

La co-ocurrencia de TEA y AACC presenta una fenomenología clínica y educativa
específica que no es simplemente la suma de ambas condiciones
\parencite{burgersveltmeijer2011}. Los individuos 2e con TEA exhiben:

\begin{itemize}[leftmargin=*, itemsep=0pt]
    \item Un \textbf{perfil cognitivo en meseta} (\textit{spiky profile}): puntuaciones
    muy elevadas en razonamiento abstracto, reconocimiento de patrones y memoria
    semántica, combinadas con puntuaciones bajas en velocidad de procesamiento y
    memoria de trabajo \parencite{assouline2012}.
    \item \textbf{Enmascaramiento (\textit{masking})} de las manifestaciones del TEA
    mediante estrategias compensatorias aprendidas, especialmente pronunciado en
    individuos con AACC y en mujeres \parencite{hull2017}.
    \item Riesgo de \textbf{diagnóstico tardío o erróneo}: las capacidades
    intelectuales superiores enmascaran las dificultades, y las dificultades
    enmascaran las capacidades, generando lo que se ha descrito como doble
    enmascaramiento \parencite{foleynicpon2017}.
    \item Presencia de \textbf{comorbilidades} frecuentes: TDAH (50--70\,\% según
    algunas series), trastornos de ansiedad, disfunciones ejecutivas y
    alexitimia \parencite{deprey2011}.
\end{itemize}

\subsection{El perfil 2e en la academia y la investigación}

La literatura específica sobre adultos con TEA y AACC en entornos universitarios y
científicos es extraordinariamente escasa. Los estudios disponibles sugieren que:

\begin{enumerate}[leftmargin=*, itemsep=2pt]
    \item Los entornos STEM tienden a atraer perfiles con características autísticas
    —atención al detalle, intereses circunscritos, tolerancia a la rutina analítica—
    en mayor proporción que otros campos profesionales \parencite{baron-cohen2001}.
    \item Las instituciones académicas ofrecen simultáneamente un ambiente
    potencialmente favorable (autonomía, estructura flexible, estimulación intelectual)
    y uno potencialmente adverso (exigencias sociales intensas, networking, evaluaciones
    por pares, precariedad laboral) para las personas 2e \parencite{price2022}.
    \item El diagnóstico tardío de TEA en académicos —frecuentemente en la tercera
    o cuarta década de vida— desencadena procesos de reidentificación biográfica y
    puede tener efectos significativos sobre la trayectoria profesional
    \parencite{raymaker2020}.
\end{enumerate}

% ═══════════════════════════════════════════════════════════
\section{Metodología}
\label{sec:metodo}
% ═══════════════════════════════════════════════════════════

\subsection{Diseño del estudio}

Se realizó un análisis bibliométrico transversal retrospectivo siguiendo las pautas
metodológicas de \textcite{donthu2021} para análisis bibliométricos en ciencias
sociales. El estudio no requiere aprobación ética al trabajar exclusivamente con
metadatos de registros académicos publicados.

\subsection{Fuente de datos}

La base de datos utilizada fue \textbf{Lens.org} (\url{https://lens.org}), plataforma
de acceso abierto mantenida por Cambia (Australia) que agrega registros de PubMed,
Crossref, OpenAlex y Microsoft Academic, alcanzando más de 200 millones de registros
académicos y 120 millones de patentes \parencite{lens2024}. Lens.org fue seleccionada
por tres razones: (1) cobertura multidisciplinar superior a bases de datos
especializadas como WoS o Scopus; (2) acceso sin restricciones para exportación y
análisis; (3) disponibilidad de identificadores propios (Lens ID) que facilitan la
deduplicación entre búsquedas.

\subsection{Estrategias de búsqueda}

Se diseñaron cinco estrategias de búsqueda booleana (Tabla~\ref{tab:strings}),
ejecutadas el 5 de junio de 2026. Las búsquedas se realizaron sobre los campos
de título, resumen y palabras clave disponibles en Lens.org.

\begin{table}[H]
\centering
\caption{Strings de búsqueda utilizados en Lens.org}
\label{tab:strings}
\small
\begin{tabularx}{\textwidth}{@{}l p{3cm} X r@{}}
\toprule
\textbf{ID} & \textbf{Denominación} & \textbf{Descripción} & \textbf{N bruto} \\
\midrule
S1 & Amplio & TEA/autismo AND AACC/gifted/twice-exceptional & 3.571 \\
S2 & Adultos profesionales & S1 AND researcher/academic/university/STEM & 999 \\
S3 & Diagnóstico tardío & S1 AND late diagnosis/masking/camouflaging & 450 \\
S4 & Neurodiversidad laboral & TEA AND neurodiversity AND academic career/employment AND gifted & 120 \\
S5 & Español/Latinoamérica & TEA/autismo AND altas capacidades/doble excepcionalidad AND investigación/academia & 360 \\
\midrule
\textbf{Total bruto} & & & \textbf{5.500} \\
\bottomrule
\end{tabularx}
\end{table}

\begin{notabox}
\textbf{Nota sobre S1 y S1-base.} Los archivos exportados S1-amplio y S1-base
resultaron ser idénticos (MD5 coincidente), confirmando que corresponden a la misma
búsqueda. Solo se utilizó S1-amplio en el análisis.
\end{notabox}

\subsection{Criterios de inclusión y exclusión}

Se incluyeron todos los tipos de documentos (artículos, libros, capítulos,
\textit{preprints}, actas de conferencias) sin restricción de idioma. Se excluyeron
únicamente los registros sin título o con metadatos insuficientes para el análisis
($n=12$). No se aplicó filtro de año de publicación para preservar la perspectiva
histórica, aunque el análisis detallado se centró en el período 2000--2026.

\subsection{Proceso de deduplicación}

Los cinco corpora se unificaron en un corpus maestro aplicando deduplicación
secuencial: (1) por DOI normalizado (minúsculas, eliminación de espacios) para los
registros que disponían de este campo; (2) por Lens ID para los registros sin DOI.
Este procedimiento eliminó 1.534 duplicados (27,9\,\%), resultando en un corpus
maestro final de \textbf{3.966 registros únicos}.

\subsection{Indicadores bibliométricos}

Se calcularon los siguientes indicadores:

\begin{itemize}[leftmargin=*, itemsep=2pt]
    \item \textbf{Productividad temporal:} distribución anual de publicaciones,
    tasa de crecimiento anual compuesto (CAGR), identificación de puntos de inflexión.
    \item \textbf{Ley de Bradford:} distribución de revistas en tres zonas de
    concentración (núcleo, zona 2, periferia) para el subconjunto de artículos.
    \item \textbf{Ley de Lotka:} distribución de productividad de autores
    (relación log-log entre número de autores y número de publicaciones).
    \item \textbf{Análisis de coautoría:} construcción de red de colaboración
    con cálculo de centralidad de grado y \textit{betweenness}.
    \item \textbf{Análisis léxico de títulos:} frecuencia de términos, comparación
    temporal pre/post 2020, \textit{WordCloud}.
    \item \textbf{Solapamiento entre estrategias:} matriz de intersección de
    Lens IDs entre los cinco corpora.
\end{itemize}

\subsection{Herramientas de análisis}

El análisis se realizó íntegramente en Python 3.12 mediante un cuaderno
Jupyter. Las bibliotecas principales fueron: \texttt{bibtexparser} (v1.4.4,
importación de archivos BibTeX), \texttt{pandas} (transformación y análisis),
\texttt{matplotlib/seaborn} (visualización), \texttt{networkx} (análisis de redes),
y \texttt{wordcloud} (visualización de frecuencias léxicas). El cuaderno y el
corpus maestro en CSV están disponibles como material suplementario.

% ═══════════════════════════════════════════════════════════
\section{Resultados}
\label{sec:resultados}
% ═══════════════════════════════════════════════════════════

\subsection{Estructura y composición del corpus}

El corpus maestro final comprende \textbf{3.966 registros únicos} distribuidos en
seis tipos de documento (Figura~\ref{fig:overview}). Los artículos de revista
constituyen la categoría dominante (71,6\,\%, $n=2.838$), seguidos de
\textit{preprints} y documentos de tipo \textit{misc} (14,7\,\%, $n=582$), capítulos
de libro (5,3\,\%, $n=212$) y libros (6,3\,\%, $n=251$). Las actas de conferencias
representan el 2,0\,\% ($n=81$).

\begin{figure}[H]
    \centering
    \includegraphics[width=\textwidth]{fig_01_corpus_overview.png}
    \caption{Estructura del corpus bibliométrico. \textit{Izquierda:} distribución
    de registros por string de búsqueda. \textit{Derecha:} composición por tipo de
    documento en el corpus maestro deduplicado.}
    \label{fig:overview}
\end{figure}

El solapamiento entre estrategias de búsqueda (Figura~\ref{fig:heatmap}) revela
patrones relevantes: S3 (diagnóstico tardío) está completamente contenido en S1
(solapamiento 100\,\%), y S2 (adultos profesionales) comparte el 95,3\,\% de sus
registros con S1. Por el contrario, S5 (español/latinoamérica) presenta un
solapamiento de apenas el 5,0\,\% con S1 y nulo con S4, lo que confirma que la
literatura en español constituye un corpus casi completamente independiente del
anglosajón — un hallazgo de notable relevancia metodológica.

\begin{figure}[H]
    \centering
    \includegraphics[width=0.75\textwidth]{fig_10_overlap_heatmap.png}
    \caption{Matriz de solapamiento entre las cinco estrategias de búsqueda,
    expresada en número de Lens IDs compartidos. La diagonal principal corresponde
    al total de cada string. El escaso solapamiento de S5 con el resto confirma la
    segmentación lingüística del campo.}
    \label{fig:heatmap}
\end{figure}

\subsection{Evolución temporal de la producción científica}

El análisis temporal revela una \textbf{trayectoria exponencial sostenida}
desde el año 2000 (Figura~\ref{fig:temporal}). En el período 2000--2010, la tasa
de crecimiento anual compuesto fue del 8,5\,\%; en el período 2010--2020 se aceleró
hasta el 16,4\,\%, y se mantuvo en el 16,7\,\% entre 2015 y 2022.

\begin{figure}[H]
    \centering
    \includegraphics[width=\textwidth]{fig_02_temporal_evolution.png}
    \caption{Evolución temporal de la producción científica. \textit{Panel superior:}
    corpus maestro con media móvil de tres años (línea roja). Se señalan dos hitos
    clave: publicación del DSM-5 (2013) y comienzo de la pandemia COVID-19 (2020).
    \textit{Panel inferior:} tendencias por string de búsqueda (2005--2026).}
    \label{fig:temporal}
\end{figure}

El año 2021 registró el volumen máximo histórico de publicaciones ($n=425$),
inmediatamente seguido por 2022 ($n=424$). El 49,1\,\% del corpus total se publicó
en los cinco años posteriores a 2020, y el 78,3\,\% corresponde a la última
década (post-2015). La caída observada en 2023 ($n=187$) y la recuperación
parcial en 2025 ($n=177$) pueden estar relacionadas con rezagos de indexación
en Lens.org para los años más recientes, más que con una contracción real del campo.

Dos hitos disciplinares parecen haber operado como catalizadores del crecimiento:
(1) la publicación del \textbf{DSM-5} en 2013, que al unificar el espectro y
eliminar el diagnóstico de síndrome de Asperger como categoría separada, estimuló
una oleada de investigación sobre diagnóstico diferencial y perfiles de alta
funcionalidad; y (2) la \textbf{pandemia de COVID-19} (2020), que puede haber
favorecido el autodiagnóstico y la búsqueda de evaluación entre adultos con perfiles
neurodivergentes durante los períodos de confinamiento.

La evolución por string (panel inferior, Figura~\ref{fig:temporal}) muestra que S1
concentra el grueso del crecimiento, mientras que S4 (neurodiversidad en entornos
laborales) y S5 (español) presentan volúmenes muy inferiores. Esto señala que,
aunque el campo general TEA--AACC ha crecido vigorosamente, las subáreas de adultos
en academia y de producción iberoamericana no han seguido el mismo ritmo.

\subsection{Análisis de tipos de documento por período}

La Figura~\ref{fig:doctypes} muestra la composición del corpus por tipo de
documento entre 2010 y 2026. Los artículos de revista mantienen una posición
dominante y estable, aunque se observa un incremento notable de
\textit{preprints/misc} a partir de 2019, consistente con la generalización
de plataformas como Zenodo, bioRxiv y PsyArXiv. La proporción de capítulos y
libros disminuye relativamente, lo que refleja el desplazamiento de la
comunicación científica hacia formatos de acceso abierto.

\begin{figure}[H]
    \centering
    \includegraphics[width=\textwidth]{fig_05_doctypes_time.png}
    \caption{Tipos de documentos por año (2010--2026). Los artículos de revista
    (azul) dominan en todos los períodos. El incremento de \textit{misc/preprints}
    (morado) desde 2019 refleja la expansión de repositorios abiertos.}
    \label{fig:doctypes}
\end{figure}

\subsection{Revistas nucleares: distribución Bradford}

El análisis de Bradford (Figura~\ref{fig:bradford}) se aplicó al subconjunto de
artículos de revista ($n=2.838$) distribuidos en \textbf{1.635 fuentes distintas}.
El núcleo Bradford (primer tercio de la producción acumulada) comprende
\textbf{113 revistas}, la zona 2 abarca 565 fuentes, y la periferia engloba
957 revistas. Esta distribución altamente dispersa (relación núcleo:total $\approx$ 1:14)
es característica de campos de investigación jóvenes e interdisciplinares
\parencite{bornmann2012}.

\begin{figure}[H]
    \centering
    \includegraphics[width=\textwidth]{fig_03_journals_bradford.png}
    \caption{Distribución Bradford. \textit{Izquierda:} top 20 revistas por
    número de artículos. \textit{Derecha:} curva acumulada en escala semilogarítmica
    con identificación de las tres zonas Bradford.}
    \label{fig:bradford}
\end{figure}

Las diez revistas más productivas del núcleo Bradford se presentan en la
Tabla~\ref{tab:revistas}. \textit{Frontiers in Psychology} lidera con 65 artículos
(2,3\,\% de todos los artículos), seguida del \textit{Journal of Autism and
Developmental Disorders} (JADD, $n=41$), referente histórico del campo, y
\textit{Focus on Exceptional Children} ($n=33$), especializada en educación
diferenciada. La presencia de \textit{Zenodo} como quinta fuente más productiva
($n=28$) confirma el peso creciente de los repositorios abiertos como canales
de difusión.

\begin{table}[H]
\centering
\caption{Top 10 revistas del núcleo Bradford (artículos de revista)}
\label{tab:revistas}
\small
\begin{tabularx}{\textwidth}{@{}l l r r@{}}
\toprule
\textbf{Rango} & \textbf{Revista} & \textbf{N} & \textbf{\% acum.} \\
\midrule
1  & Frontiers in Psychology                              & 65  & 2.3\,\% \\
2  & Journal of Autism and Developmental Disorders        & 41  & 3.7\,\% \\
3  & Focus on Exceptional Children                        & 33  & 4.9\,\% \\
4  & PLOS ONE                                             & 34  & 6.1\,\% \\
5  & Zenodo (CERN -- repositorio abierto)                 & 28  & 7.1\,\% \\
6  & Frontiers in Psychiatry                              & 26  & 8.0\,\% \\
7  & Education Sciences                                   & 20  & 8.7\,\% \\
8  & Scientific Reports                                   & 20  & 9.4\,\% \\
9  & Autism: The International Journal                    & 18  & 10.0\,\% \\
10 & Frontiers in Human Neuroscience                      & 14  & 10.5\,\% \\
\bottomrule
\end{tabularx}
\end{table}

La heterogeneidad de las revistas del núcleo —que abarca desde psicología clínica
(JADD, \textit{Autism Research}) hasta educación especial (\textit{Focus on
Exceptional Children}, \textit{Journal for the Education of Gifted Young Scientists})
pasando por ciencias de la salud (\textit{Frontiers in Psychiatry}) y revistas
multidisciplinares (PLOS ONE)— evidencia la naturaleza \textbf{transversal e
interdisciplinar} del campo TEA--AACC.

\subsection{Análisis de autoría: distribución de Lotka}

El corpus presenta 20.986 autores únicos (Figura~\ref{fig:lotka}). La distribución
de productividad sigue el patrón predicho por la Ley de Lotka: el 90,9\,\% de los
autores tiene una única publicación en el corpus, frente al 62\,\% teórico de Lotka
para coeficiente $n=2$. Esta desviación hacia una concentración aún mayor sugiere
que el campo está en proceso de consolidación, con una base amplísima de
contribuyentes ocasionales y un núcleo muy reducido de investigadores especializados.

\begin{figure}[H]
    \centering
    \includegraphics[width=\textwidth]{fig_04_authors_lotka.png}
    \caption{Distribución de productividad de autores (Ley de Lotka). \textit{Izquierda:}
    top 20 autores por número de publicaciones. \textit{Derecha:} distribución log-log
    de frecuencias de productividad.}
    \label{fig:lotka}
\end{figure}

Solo 24 autores tienen 10 o más publicaciones en el corpus, y 130 autores tienen
5 o más. Los autores más productivos (Tabla~\ref{tab:autores}) incluyen a
\textbf{R. Plomin} ($n=55$), genetista conductual del King's College London cuya
producción sobre heredabilidad de la inteligencia y el comportamiento autístico
es consistente con el foco del corpus; \textbf{Megan Foley-Nicpon} ($n=19$) y
\textbf{Susan G. Assouline} ($n=17$) de la Universidad de Iowa, quienes constituyen
el equipo referente mundial en investigación empírica sobre 2eASD; y
\textbf{Athanasios Drigas} ($n=13$), investigador griego especializado en tecnologías
educativas para personas con necesidades especiales.

\begin{table}[H]
\centering
\caption{Top 15 autores más productivos del corpus maestro}
\label{tab:autores}
\small
\begin{tabularx}{\textwidth}{@{}r l r l@{}}
\toprule
\textbf{Rango} & \textbf{Autor} & \textbf{N pubs.} & \textbf{Afiliación/especialidad conocida} \\
\midrule
1  & R. Plomin            & 55 & King's College London -- Genética conductual \\
2  & Lisa Corwin          & 22 & Colorado -- Educación STEM \\
3  & T. Chalder           & 21 & King's College London -- Psicología clínica \\
4  & Megan Foley-Nicpon   & 19 & Univ. Iowa -- 2e ASD especialista \\
5  & Susan G. Assouline   & 17 & Univ. Iowa -- 2e ASD especialista \\
6  & S. Wessely           & 16 & King's College London -- Psiquiatría \\
7  & Meredith Henry       & 14 & STEM education \\
8  & Shayla Shorter       & 14 & STEM education \\
9  & Louise Charkoudian   & 14 & STEM education \\
10 & Jennifer Heemstra    & 14 & STEM education \\
11 & Athanasios Drigas    & 13 & NCSR Demokritos (Grecia) -- TIC y NEE \\
12 & R. Moss-Morris       & 13 & King's College London \\
13 & J. DeFries           & 13 & Univ. Colorado -- Genética \\
14 & P. Dale              & 12 & Genética del lenguaje \\
15 & S. Petrill           & 12 & Ohio State -- Genética conductual \\
\bottomrule
\end{tabularx}
\end{table}

\begin{notabox}
\textbf{Observación.} La alta concentración de publicaciones en el grupo Plomin
(n=55), Chalder, Wessely y Moss-Morris (todos King's College) y el clúster de
educación STEM (Corwin, Henry, Shorter, Heemstra) sugiere que \textit{parte del
corpus S1 puede incluir registros que comparten autores pero no son específicamente
sobre TEA+AACC} (p. ej., estudios sobre fatiga crónica de Wessely). Este es un
hallazgo relevante sobre los límites del string S1 y justifica la importancia de
los strings más específicos S2--S4 para el nicho de interés de este estudio.
\end{notabox}

\subsection{Red de co-autoría}

La red de co-autoría (Figura~\ref{fig:red}) construida con autores de $\geq3$
publicaciones revela una topología de \textbf{red de escala libre} con un núcleo
denso central y numerosos autores periféricos con pocas conexiones. Los nodos con
mayor betweenness centrality (autores con mayor capacidad de mediación entre
subcomunidades) incluyen a Foley-Nicpon, Nicolson y Hulme, lo que confirma su
papel de puentes entre las comunidades de investigación sobre TEA, sobre AACC y
sobre psicología educativa.

\begin{figure}[H]
    \centering
    \includegraphics[width=0.9\textwidth]{fig_07_coauthorship_network.png}
    \caption{Red de co-autoría. Nodos: autores con $\geq3$ publicaciones. Tamaño
    del nodo: proporcional al número de publicaciones. Color: betweenness centrality
    (amarillo~$\rightarrow$~rojo = mayor centralidad). Las etiquetas corresponden
    a los 15 autores con mayor betweenness.}
    \label{fig:red}
\end{figure}

\subsection{Análisis léxico de títulos}

El análisis de frecuencia de términos en títulos del corpus maestro
(Figura~\ref{fig:wordcloud}) muestra que los términos más frecuentes son
\textit{autism} (537), \textit{children} (501) y \textit{education} (413).
La presencia masiva de \textit{children} (501 ocurrencias) frente a
\textit{adults} (ausente entre los 40 primeros) refleja cuantitativamente
el sesgo del campo hacia la población infantil.

\begin{figure}[H]
    \centering
    \includegraphics[width=\textwidth]{fig_06_wordcloud_terms.png}
    \caption{Análisis léxico de títulos del corpus maestro.
    \textit{Izquierda:} nube de palabras proporcional a la frecuencia.
    \textit{Derecha:} frecuencia absoluta de los 30 términos más comunes.
    Nótese la dominancia de \textit{children}, \textit{students} y
    \textit{school} frente a la escasez de términos relacionados con adultos
    o contextos laborales.}
    \label{fig:wordcloud}
\end{figure}

El análisis comparativo pre/post 2020 reveló los términos con mayor crecimiento
relativo en la producción reciente, entre los que destacan \textit{neurodiversity},
\textit{twice-exceptional}, \textit{masking} y \textit{higher education}, lo que
sugiere un incipiente giro temático hacia los intereses de este estudio.

\subsection{Literatura en español y contexto latinoamericano}

El string S5 recuperó 360 registros en español (5,0\,\% de solapamiento con el
corpus anglosajón), de los cuales 343 permanecieron en el corpus maestro.
La evolución temporal (Figura~\ref{fig:spanish}) muestra un crecimiento desde
prácticamente cero antes de 2012 hasta un pico de 47 publicaciones en 2022.
Las revistas más productivas son el \textit{International Journal of Developmental
and Educational Psychology -- INFAD} ($n=5$), el \textit{Electronic Journal of
Research in Education Psychology} ($n=4$), y la \textit{Revista de Psiquiatría
Infanto-Juvenil} ($n=4$).

\begin{figure}[H]
    \centering
    \includegraphics[width=\textwidth]{fig_08_spanish_literature.png}
    \caption{Literatura iberoamericana. \textit{Izquierda:} evolución temporal
    (2000--2026). \textit{Derecha:} top 10 revistas en español con mayor
    producción sobre TEA + AACC.}
    \label{fig:spanish}
\end{figure}

El análisis léxico de los títulos en español (Figura~\ref{fig:wc_es}) muestra
dominancia de \textit{educación}, \textit{inclusión} y \textit{necesidades},
con presencia de \textit{altas capacidades} y \textit{espectro}, pero notable
ausencia de términos relacionados con adultos profesionales, academia o
investigación. El término \textit{ensino} (del portugués) refleja la
contribución significativa de Brasil al corpus iberoamericano.

\begin{figure}[H]
    \centering
    \includegraphics[width=\textwidth]{fig_09_spanish_wordcloud.png}
    \caption{Análisis léxico de títulos en español/portugués (S5).
    La terminología refleja un enfoque predominantemente escolar e
    inclusión educativa, con escasa presencia de términos asociados
    a adultos o contextos universitarios.}
    \label{fig:wc_es}
\end{figure}

\subsection{El contexto colombiano: un caso de invisibilidad científica}

La búsqueda de registros con mención explícita a Colombia identificó únicamente
\textbf{8 documentos} en todo el corpus (Tabla~\ref{tab:colombia}), todos ellos
concentrados en los años 2020--2025, y ninguno centrado específicamente en la
intersección TEA--AACC en adultos o investigadores. Seis de los ocho registros
aparecen en la \textit{Revista Colombiana de Educación} y revistas afines, y los
temas abordados oscilan entre políticas de educación inclusiva general y creencias
del profesorado sobre el autismo.

\begin{table}[H]
\centering
\caption{Registros con referencia al contexto colombiano identificados en el corpus}
\label{tab:colombia}
\small
\begin{tabularx}{\textwidth}{@{}r l p{7.5cm} l@{}}
\toprule
\textbf{Año} & \textbf{Fuente} & \textbf{Título (fragmento)} & \textbf{Foco} \\
\midrule
2020 & Rev. Colombiana de Educación & Aproximaciones institucionales a estudiantes con discapacidad & Ed. inclusiva \\
2020 & Rev. Colombiana de Educación & Educación de sujetos con discapacidad en Colombia & Historia \\
2020 & Magis RIIE & Superdotación e identidad: lectura desde Colombia & AACC \\
2020 & Act. Invest. Educación & Legislación colombiana y evaluación escolar NEE & Política \\
2020 & -- & Investigación en Doctorado en Ciencias de la Educación & Ed. superior \\
2022 & Rev. Colombiana de Educación & Videojuegos y atención a la diversidad & TIC \\
2022 & -- & Creencias sobre TEA en maestras de educación inicial & TEA \\
2025 & Rev. Lat. Ciencias Sociales & Formación de personas con altas capacidades en Rionegro & AACC \\
\bottomrule
\end{tabularx}
\end{table}

Este resultado cuantifica de manera objetiva lo que los investigadores del campo
conocen cualitativamente: \textbf{Colombia carece de una comunidad investigadora
consolidada} sobre TEA--AACC, y prácticamente no existen estudios sobre adultos
con doble excepcionalidad en entornos académicos o científicos colombianos. La
única referencia cercana al foco de este estudio es el artículo de \textit{Magis}
sobre superdotación e identidad desde una perspectiva colombiana (2020), que aborda
la AACC pero no específicamente su co-ocurrencia con TEA en adultos profesionales.

% ═══════════════════════════════════════════════════════════
\section{Discusión}
\label{sec:discusion}
% ═══════════════════════════════════════════════════════════

\subsection{Fortalezas del campo: crecimiento y consolidación temática}

Los resultados confirman que la investigación sobre TEA--AACC ha alcanzado un
volumen y una tasa de crecimiento que justifican hablar de campo emergente en
proceso de consolidación. El CAGR del 16,4\,\% en la última década supera con
creces el crecimiento de la producción científica general (estimado en 4--5\,\%
anual \parencite{bornmann2021}) y es comparable al de otras áreas de alta
actividad como la investigación sobre microbioma \parencite{chen2022}.

La aparición de revistas especializadas de alto impacto en el núcleo Bradford
—JADD, \textit{Autism}, \textit{Gifted Child Quarterly}— y la presencia de equipos
de investigación consolidados (Foley-Nicpon/Assouline en Iowa; Drigas en Atenas;
Reis/Renzulli en Connecticut) indican que existe una infraestructura disciplinar
mínima viable. Los instrumentos de evaluación están mejorando: la adaptación del
ADOS y el ADI-R para perfiles de alta capacidad \parencite{foleynicpon2017}, y el
desarrollo de baterías específicamente diseñadas para la evaluación de doble
excepcionalidad \parencite{cederberg2018}, representan avances metodológicos
significativos.

\subsection{Debilidades estructurales: el sesgo hacia la infancia y la escuela}

La debilidad más palmaria del campo es el \textbf{sesgo pediátrico y escolar}.
El término \textit{children} aparece en 501 títulos frente a una presencia marginal
de \textit{adults}, \textit{university} o \textit{researcher}. Solo el 6,5\,\% del
corpus menciona explícitamente adultos o entornos de educación superior. Este sesgo
no es neutral: refleja el supuesto implícito de que la doble excepcionalidad es
un fenómeno de la infancia que se ``resuelve'' en la adultez, cuando la evidencia
clínica indica lo contrario \parencite{raymaker2020}.

La lógica de intervención temprana domina el imaginario del campo: identificar,
diagnosticar e intervenir en los primeros años de escolarización para optimizar
las trayectorias educativas. Esta lógica, aunque legítima, invisibiliza a los
adultos que llegaron a la vida profesional sin diagnóstico, navegando sus carreras
académicas o científicas en condiciones de mayor dificultad y con acceso limitado
a apoyos institucionales \parencite{price2022}.

\subsection{El enmascaramiento como factor explicativo de la invisibilidad}

El \textit{masking} o enmascaramiento —la estrategia de suprimir o disimular
conductas autísticas para encajar en contextos neurotípicos— es especialmente
pronunciado en individuos con AACC, que disponen de mayores recursos cognitivos
para la imitación social y la compensación conductual \parencite{hull2017}.
En contextos académicos y científicos, donde la presión de conformidad social
es alta y la evaluación del rendimiento relativamente objetiva, el enmascaramiento
puede resultar especialmente eficaz \textit{y} especialmente costoso: el coste
neurológico del enmascaramiento sostenido se asocia con burnout, disociación y
crisis de salud mental \parencite{mandy2019}.

La casi total ausencia de literatura sobre masking en el corpus —solo 1 artículo
con el término en el título— frente a la abundante evidencia clínica de su
prevalencia en adultos de alto rendimiento, constituye una brecha de investigación
de primera magnitud.

\subsection{Debilidades metodológicas de los estudios existentes}

Las revisiones disponibles \parencite{gelbar2022, luor2021} han documentado
sistemáticamente las debilidades metodológicas del campo: (a) tamaños muestrales
pequeños, frecuentemente $n < 30$; (b) ausencia de grupos de control adecuados;
(c) falta de homogeneidad en los criterios de definición tanto de AACC como de TEA;
(d) predominio de metodologías cualitativas y estudios de caso; y (e) reclutamiento
por conveniencia en poblaciones clínicas, que introduce sesgo de selección. A estos
problemas debe añadirse la escasez de estudios longitudinales que permitan seguir
trayectorias desde la infancia hasta la vida adulta profesional.

\subsection{La perspectiva latinoamericana: entre el rezago y la oportunidad}

La producción iberoamericana sobre TEA--AACC representa el 8,6\,\% del corpus
global con apenas un 5,0\,\% de solapamiento con la literatura anglosajona.
Esto significa que ambas comunidades científicas operan prácticamente en mundos
paralelos, con escasa transferencia de conocimiento bidireccional.

Las particularidades del contexto latinoamericano justifican una agenda investigadora
propia: (a) la prevalencia estimada de TEA en la región es inferior a la reportada
en países de altos ingresos \parencite{salari2022}, probablemente por infradiagnóstico
más que por diferencias reales de prevalencia; (b) los sistemas de educación
especial tienen menor capacidad de identificación temprana; (c) la definición
y los criterios para el reconocimiento de AACC varían enormemente entre países; y
(d) el estigma asociado al diagnóstico de TEA es percibido como mayor en contextos
culturales latinoamericanos \parencite{conejeros2018}.

\subsection{Colombia: brecha crítica y agenda urgente}

Los ocho documentos colombianos identificados confirman que Colombia no tiene
prácticamente investigación sobre la intersección TEA--AACC. Las implicaciones
son preocupantes a nivel de política educativa y sanitaria: sin datos propios,
las decisiones de política se basan en evidencia extranjera de dudosa transferibilidad
cultural; sin investigación sobre adultos y académicos, el sistema universitario
colombiano carece de bases para diseñar apoyos institucionales para su personal
docente e investigador neurodivergente; y sin visibilidad científica, la comunidad
de profesionales con TEA y AACC en Colombia permanece sin reconocimiento y sin voz
en los espacios de decisión.

Colombia posee, sin embargo, condiciones estructurales para impulsar esta agenda:
una comunidad investigadora en educación especial y psicología educativa en
crecimiento (visibles en la \textit{Revista Colombiana de Educación} y en redes
como RUDECOLOMBIA), una legislación de discapacidad e inclusión relativamente
avanzada en el contexto regional (Ley 1618 de 2012), y un sistema de ciencia y
tecnología —MINCIENCIAS— que ha comenzado a financiar investigación sobre
neurodiversidad.

% ═══════════════════════════════════════════════════════════
\section{Limitaciones del estudio}
\label{sec:limitaciones}
% ═══════════════════════════════════════════════════════════

Este estudio presenta varias limitaciones que deben tenerse en cuenta al
interpretar los resultados:

\begin{enumerate}[leftmargin=*, itemsep=3pt]
    \item \textbf{Ausencia de datos de citación.} El formato BibTeX exportado por
    Lens.org no incluye recuentos de citas, lo que impide calcular el índice~$h$
    del campo ni el impacto ponderado por citas de revistas y autores.

    \item \textbf{Rezago de indexación.} Lens.org puede presentar retrasos de
    indexación para los años más recientes (2023--2026), lo que explicaría
    parcialmente la caída de producción observada en 2023.

    \item \textbf{Análisis léxico limitado a títulos.} Por ausencia del campo
    \textit{keywords} en los registros exportados, el análisis terminológico se
    restringió a títulos, que no siempre reflejan fielmente el contenido.

    \item \textbf{Ausencia de análisis de co-citación.} Sin datos de referencias,
    no fue posible construir redes de co-citación o de acoplamiento bibliográfico,
    que habrían permitido mapear con mayor precisión los clústeres temáticos del campo.

    \item \textbf{Cobertura del contexto colombiano.} La búsqueda de registros
    colombianos se basó en términos léxicos en títulos y nombres de revistas, lo
    que puede haber omitido trabajos de autores colombianos publicados en revistas
    internacionales sin mención explícita al país.

    \item \textbf{Strings de búsqueda.} Aunque los strings fueron diseñados con
    criterio experto, la amplia inclusividad de S1 puede haber incorporado registros
    que no son específicamente sobre TEA+AACC (como se señaló para el clúster Plomin-Wessely).
\end{enumerate}

% ═══════════════════════════════════════════════════════════
\section{Líneas prioritarias de investigación futura}
\label{sec:futuro}
% ═══════════════════════════════════════════════════════════

A partir de los gaps identificados en este análisis, se proponen seis líneas de
investigación con alta prioridad:

\begin{enumerate}[leftmargin=*, label=\textbf{LP\arabic*.}, itemsep=4pt]
    \item \textbf{Estudios de prevalencia y caracterización en adultos académicos.}
    Investigación epidemiológica sobre la prevalencia del perfil 2e (TEA+AACC) en
    población investigadora y académica universitaria, con atención específica a la
    tasa de diagnóstico tardío y las consecuencias sobre la trayectoria profesional.

    \item \textbf{Enmascaramiento en contextos científicos.}
    Estudios cualitativos y mixtos sobre las estrategias de masking utilizadas por
    investigadores con TEA, sus costes neurológicos y psicológicos, y su relación
    con el síndrome de impostora/impostor y el burnout académico.

    \item \textbf{Interseccionalidad género, TEA y AACC en academia.}
    Las mujeres autistas con altas capacidades constituyen un subgrupo con doble
    enmascaramiento (autístico y de género) particularmente vulnerable e
    investigado. Se necesitan estudios con diseños sensibles al género y al
    ciclo vital.

    \item \textbf{Perspectiva latinoamericana comparada.}
    Estudios comparativos entre países de la región sobre diagnóstico, políticas
    de apoyo y trayectorias profesionales de adultos 2e, con metodologías adaptadas
    al contexto cultural latinoamericano.

    \item \textbf{Colombia: investigación de base.}
    Estudios exploratorios sobre el perfil 2e en universidades e institutos de
    investigación colombianos, incluyendo: caracterización de necesidades de apoyo
    institucional, análisis de barreras diagnósticas, y mapeo de buenas prácticas
    en instituciones de educación superior.

    \item \textbf{Diseño de intervenciones específicas para adultos 2e en academia.}
    Desarrollo y validación de programas de apoyo institucional (mentoring, ajustes
    razonables, formación de pares) adaptados a las necesidades específicas de
    investigadores y docentes con TEA y AACC.
\end{enumerate}

% ═══════════════════════════════════════════════════════════
\section{Conclusiones}
\label{sec:conclusiones}
% ═══════════════════════════════════════════════════════════

Este análisis bibliométrico, basado en un corpus de 3.966 registros recuperados de
Lens.org y procesados íntegramente con herramientas de Python, ofrece la
caracterización más actualizada y de mayor cobertura disponible sobre la producción
científica en la intersección TEA--AACC. Los hallazgos principales son:

\begin{itemize}[leftmargin=*, itemsep=3pt]
    \item El campo ha experimentado un crecimiento exponencial sostenido (CAGR
    16,4\,\% en 2010--2020), con el 78,3\,\% del corpus publicado en la última
    década, lo que evidencia un campo emergente en aceleración.

    \item La producción está concentrada en un pequeño núcleo Bradford de 113
    revistas y un número muy reducido de autores especializados (24 con $\geq10$
    publicaciones), lo que refleja la juventud e interdisciplinaridad del campo.

    \item Existe un sesgo estructural, cuantificado y estadísticamente inequívoco,
    hacia la población infantil y los entornos escolares, con el 93,5\,\% del corpus
    sin mención de adultos o academia en los títulos.

    \item La literatura en español y la producción latinoamericana constituyen
    un corpus casi independiente del anglosajón (5\,\% de solapamiento), con
    una agenda temática dominada por la educación inclusiva escolar y prácticamente
    sin investigación sobre adultos profesionales.

    \item Colombia tiene una presencia científica mínima en el campo (8 registros
    en todo el corpus global), confirmando la necesidad urgente de una agenda
    investigadora propia.
\end{itemize}

La doble excepcionalidad TEA--AACC en investigadores y académicos no es una
rareza clínica ni un problema de nicho: es una realidad presente y silenciada en
las instituciones académicas de todo el mundo, incluyendo las colombianas. Los
sistemas universitarios que no reconozcan y apoyen activamente a su personal
neurodivergente de alta capacidad perderán una fuente de pensamiento divergente,
originalidad científica y creatividad intelectual que ninguna institución que
aspire a la excelencia puede permitirse ignorar.

% ═══════════════════════════════════════════════════════════
% BIBLIOGRAFÍA
% ═══════════════════════════════════════════════════════════

\newpage
\printbibliography[title={Referencias}]

% ═══════════════════════════════════════════════════════════
% APÉNDICES
% ═══════════════════════════════════════════════════════════

\newpage
\appendix

\section{Strings de búsqueda completos}
\label{app:strings}

\subsection*{S1 — Amplio}
\begin{verbatim}
("autism spectrum disorder" OR "ASD" OR "autism")
AND
("giftedness" OR "gifted" OR "high ability" OR "high abilities"
 OR "twice exceptional" OR "dual exceptionality"
 OR "intellectually gifted")
\end{verbatim}

\subsection*{S2 — Adultos profesionales}
\begin{verbatim}
("autism spectrum disorder" OR "ASD" OR "autistic")
AND
("giftedness" OR "high ability" OR "twice exceptional"
 OR "high intellectual potential")
AND
("researcher" OR "scientist" OR "academic" OR "university"
 OR "higher education" OR "STEM" OR "professor")
\end{verbatim}

\subsection*{S3 — Diagnóstico tardío y enmascaramiento}
\begin{verbatim}
("autism spectrum disorder" OR "ASD")
AND
("giftedness" OR "high ability" OR "intellectually gifted")
AND
("late diagnosis" OR "late identified" OR "masking"
 OR "camouflaging" OR "compensation" OR "adult diagnosis")
\end{verbatim}

\subsection*{S4 — Neurodiversidad en entornos laborales}
\begin{verbatim}
("autistic" OR "autism spectrum")
AND ("neurodiversity" OR "neurodivergent")
AND ("academic career" OR "research career" OR "workplace"
     OR "employment" OR "scientists" OR "professionals")
AND ("high ability" OR "gifted" OR "twice exceptional")
\end{verbatim}

\subsection*{S5 — Español/Latinoamérica}
\begin{verbatim}
("trastorno del espectro autista" OR "TEA" OR "autismo")
AND
("altas capacidades" OR "superdotación" OR "doble excepcionalidad"
 OR "alta capacidad intelectual")
AND
("investigación" OR "academia" OR "universidad" OR "investigadores")
\end{verbatim}

\section{Estadísticas del corpus maestro}
\label{app:stats}

\begin{table}[H]
\centering
\caption{Estadísticas completas del corpus maestro}
\label{tab:stats}
\small
\begin{tabular}{@{}lr@{}}
\toprule
\textbf{Indicador} & \textbf{Valor} \\
\midrule
Registros totales (bruto, 5 strings)    & 5.500 \\
Duplicados eliminados                    & 1.534 (27,9\,\%) \\
Corpus maestro final                     & 3.966 \\
\quad Artículos de revista              & 2.838 (71,6\,\%) \\
\quad Misc/Preprints                    & 582 (14,7\,\%) \\
\quad Libros                            & 251 (6,3\,\%) \\
\quad Capítulos                         & 212 (5,3\,\%) \\
\quad Actas de conferencias             & 81 (2,0\,\%) \\
\quad Reportes técnicos                 & 2 (0,1\,\%) \\
\midrule
Revistas únicas (artículos)             & 1.635 \\
Zona 1 Bradford (núcleo)                & 113 revistas \\
Zona 2 Bradford                         & 565 revistas \\
Zona 3 Bradford (periferia)             & 957 revistas \\
\midrule
Autores únicos mencionados              & 20.986 \\
Autores con 1 sola publicación          & 19.078 (90,9\,\%) \\
Autores con $\geq$5 publicaciones       & 130 (0,6\,\%) \\
Autores con $\geq$10 publicaciones      & 24 (0,1\,\%) \\
\midrule
Rango temporal del corpus               & 1834--2026 \\
Año de mayor producción                 & 2021 ($n=425$) \\
CAGR 2000--2010                         & 8,5\,\% \\
CAGR 2010--2020                         & 16,4\,\% \\
Registros post-2015                     & 78,3\,\% \\
Registros post-2020                     & 49,1\,\% \\
\midrule
Registros en español (S5)               & 343 (8,6\,\%) \\
Solapamiento S5--S1                     & 5,0\,\% \\
Registros con referencia a Colombia     & 8 (0,2\,\%) \\
\bottomrule
\end{tabular}
\end{table}

\section{Material suplementario}
\label{app:suplementario}

El material suplementario de este artículo incluye:

\begin{itemize}
    \item \texttt{bibliometria\_TEA\_AACC.ipynb} — Cuaderno Jupyter con todo el
    código de análisis, reproducible con Python~3.10+.
    \item \texttt{corpus\_maestro\_TEA\_AACC\_academia.csv} — Corpus maestro
    deduplicado en formato CSV (UTF-8).
    \item \texttt{fig\_01} a \texttt{fig\_10} — Figuras del análisis en formato PNG.
\end{itemize}

Los archivos BibTeX originales (S1--S5) fueron exportados de Lens.org el
5 de junio de 2026 y están disponibles bajo solicitud a los autores.

\end{document}
